%TCIDATA{Version=5.50.0.2890}
%TCIDATA{LaTeXparent=0,0,NoticeDKE.tex}


The determination of the current density profile is a crucial issue for next
step tokamaks like ITER, since it may directly modify plasma performances.
Up to now, most of the tools where designed for simplified magnetic
configuration, or using crude physical assumptions that consequently reduce
considerably the confidence in the results. Accurate and realistic kinetic
calculations for the electron population in magnetized plasmas where
therefore an important step towards a consistent modeling of the current
drive process.

For this purpose, a completely new numerical solver has been designed which
contains most of the salient features of the physics in hot plasmas, where
the weak collision regime holds: relativistic collision operator,
bounce-averaging for trapped and passing particles, arbitrary magnetic
configuration, first order neoclassical corrections. A special attention has
been paid to derive in a rigorous and consistent manner all the equations,
in particular for the interplay between inductive, non-inductive and
bootstrap currents. The conservative nature of the equations has been the
guideline in this work, especially to have a clear derivation of radial
transport equation for the fast electron population in arbitrary magnetic
configuration when bounce averaging is concerned.

The numerical part has been derived with the same spirit, keeping a
systematic approach for all types of operators in the discretization
technique. This point is especially important for a consistent evaluation of
the current drive efficiency, a key parameter in order to evaluate
performances of the method used for driving current. In particular, the
formalism is general so that arbitrary type of mechanismes may be
incorporated for possible synergistic effect. As for the analytical part, a
detail analysis of the conservative nature of the numerical code has been
performed. A clear and comprehensive discussion of all the aspects is
presented, which is crucial for a reliable numerical tool. It is shown that
the both for the zero- and first-order equations, the numerical solver is
naturally conservative, in the sense that no need of arbitrary external
source or normalization factor is necessary in order to compensate possible
hidden numerical particle leaks which are the consequence of an incorrect
discrete projection of the equations on the numerical grids. This reference
work is critical in order to make the code transparent for possible further
evolutions.

Advanced numerical matrix factorization techniques have shown that
considerable savings may be expected, not only for storage purposes but also
for the time duration needed in determining the steady-state current drive
solution on the collision and fast electron transport time scales. The
pionneering work performed initially for the $2-D$ problem in velocity space
in a fully implicit manner (reverse time scheme that is inconditionnaly
stable with respect to the integration time step)(\cite{sho94},\cite{shk97}
and \cite{pey98}) has been fully extended to the $3-D$ case (\cite
{pey03c1_RF15}), including the complexity arising from the radial dependence
of the trapped-passing boundary for bounce-averaged equations. With the
incomplete $LU$ technique (\cite{saa03}), it has been possible to converge
in few iterations towards the steady-state current drive solution on
standard computers, taking advantage of the high sparcity of the
Fokker-Planck and electron drift kinetic equation matrices. This tool is
therefore fully designed to be incorporated in a chain of codes, for
realistic self-consistent calculations, in which plasma magnetic
equilibrium, wave propagation and absorption, energy and particle transport
must be evaluating as a function of time and space (\cite{bas03}). In is
important to notice that special attention has been paid to the delicate
problem of coupling the kinetic solver with ray-tracing. This amazing feat
has been made easily possible with the use of the compact MatLab numerical
environment for solving the huge linear system of equations (\cite{mat02}).

Furthermore, numerous moments of the electron distribution function are
evaluated like the RF power, magnetic ripple and runaway losses, fast
electron bremsstrahlung and trapped fraction. Moreover, the code may also be
used for specific physical studies which imply time evolution, and a
standard Crank-Nicholson time-scheme has been also incorporated. In the same
spirit, simplified collision models may be used, like the Lorentz operator,
for which analytical expressions may be easily derived. This point is
crucial for an effective benchmarking, even in a complex magnetic topology,
like for ITER.

At this stage, the code may be considered as mature for the initial goal
that was considered. Several additional benchmarking must be certainly done
in order to enhance the realibility on several quantities, but most of the
resultats obtained by the code are already robust. The domain of stabitility
related to the drop tolerance parameter for the approximate matrix
factorization technique needs also refined estimates for optimizing the
numerical method.

Several enhancement of the code may be easily foreseen in a near future. In
the framework, one of the most important issue is to use a correct
quasilinear diffusion operator for the wave-particle interaction. Indeed,
like for most kinetic solvers available today, the code use the well known
Kennel-Engelman-Lerche expression that is derived for plane waves in
homogenous plasmas. Regarding the assumptions, it is necessary to use a more
precise quasilinear diffusion operator that is decuded from the Hamiltonian
theory in toroidal configuration (\cite{eri94}). This requires further
analytical developments, using the appropriate coordinate system, where the
effective banana width of the particle is fully taken into account. In this
framework, the description of the wave-induced particle transport becomes
natural, and the code in its present conservative form fully allows this
kind of description that may play a fundamental role for the RF power
absorption, but also for thermal particle pinch in steady-state conditions (%
\cite{hoa03}). Furthermore, once the bounce-averaging is replaced by
trajectory-averaging, it is natural to extend the code for multiple ion
dynamics. In that case, parallel processing may be foreseen for each
species, in order to optimize calculations.

Finally, an ultimate though important evolution for the code may concern
non-linear problems. So far, the collision operator is linearized for a
Maxwellian bulk, that is a reasonable assumption in most tokamak conditions
today. However, some experimental evidences that the bulk could be
non-Maxwellian in presence of strong electron cyclotron heating and current
drive suggest that non-linear effects could take place. This may be also the
case in reactor conditions when $\alpha $-particle population is damping on
the electrons. In that case, the neoclassical theory for non-Maxwellian
distribution must be previously revisited before any numerical
implementation.

Besides the development of modern tools that is effectively described in
this document and will be useful for a detailed understanding of the current
drive in tokamak, this work has allowed to highlight the subtile interplay
between physics and numerics. This is an good example of what should be
performed for advanced realistic simulations of complex dynamical systems.
It enables to draw in a very clear manner paths for refined studies, based
on robust formalisms that takes into account in a rigorous way the physics
principles. The method here presented has therefore broad potential
applications for other challenging physics domains, like for the numerical
determination of the full-wave problem when very short wavelengths must be
considered as compared to the machine size and the gradient lengths.
